\section{専門職・継続管理・要求性能を別の軸に置く}\label{theory:compliance:management}
高度な専門能力を持つ人材の養成を目的とする専門職大学院の課程は、
博士後期課程そのものではない。専門職の学位と研究の博士課程を区別する。
法務博士という学位名の「博士」だけから博士後期課程と同一視しない。\cite{eng52_professional}

BCMは障害に備える計画を作り、訓練・点検・見直しを続けて事業継続能力を管理する活動、
BCPはその事業継続計画である。計画書を置くだけで管理の全過程が完了するわけではない。
当時の理解は平成21年版等の資料から確認し、後年の制度改訂を遡って加えない。\cite{eng52_bcp}

性能の要求は、果たすべき機能・安全・使用性等を定める。
製作仕様は材料・寸法・方法等を指定する手段であり、両者を同じ定義にしない。
性能規定で要求の満足を検証することと、仕様を守ればよいという評価を区別する。
現実の要求水準には性能と仕様が併存する場合もあるので、一方だけが常に許されるとはしない。
信頼性は必要な機能を条件・時間の中で遂行する度合いを確率的に扱う別の軸である。\cite{eng52_performance}

\subsection{VEの価値と品質の最低条件}
VEは必要な機能を確実に果たし、ライフサイクルを含む資源・コストと機能を比較して価値を改善する。
説明用の比\(V=F/C\)では、分母を下げるだけで機能の必要水準を失う変更を価値向上とはしない。
安い価格と、設計・保守・更新を含む全体コストを分ける。\cite{eng52_ve}

\subsection{内部統制の四目的と時点}
平成19年の基準では、業務の有効性・効率性、財務報告の信頼性、
事業に関わる法令等の遵守、資産の保全という四目的が示される。
組織の全員が業務の中で遂行するプロセスで、互いに関連する目的である。
経営者の整備責任と、その他の者が運用に関与することは両立する。
内部統制全体の目的と、金融商品取引法の財務報告に係る評価・報告・監査の対象を分ける。\cite{eng52_control}
平成23年度の表現を、後年の報告範囲の変更後の用語に無断置換しない。
